\section{An Explicit Retained Field and Its Hecke Factors}\label{ar:field}
The analytic estimate uses a fixed subfield of every sufficiently large
field in Theorem~\ref{tw:field-family}. We construct this subfield by
Kummer theory and express its zeta function as a product of Dirichlet and
Hecke functions. Exact local calculations then determine the conductors
and Euler polynomials needed to evaluate these factors.

\subsection{A Catalog of Norm Extensions}\label{ar:catalog}
Retain $V=\operatorname{Gal}(E/\mathbb Q)=\mathbb F_2^7$ with the
coordinates of Section~\ref{tw:tower}. For a rational squareclass $A$
supported on $-1,2,3,5,7,11,13$, write $\chi_A\in V^*$ for its Kummer
character. Define $\alpha_i=x_i+y_i\sqrt{A_i}$ by the following seventeen rows.
Each row satisfies
$N_{\mathbb Q(\sqrt{A_i})/\mathbb Q}(\alpha_i)=B_i z_i^2$.
\begin{equation}\label{ar:norm-table}
\begin{array}{r|r|r|r|r|r}
 i&A_i&B_i&x_i&y_i&z_i\\\hline
0&-1&65&4&7&1\\
1&-55&91&6&1&1\\
2&-2&33&1&4&1\\
3&-55&14&1&1&2\\
4&-55&26&7&1&2\\
5&-7&2&1&1&2\\
6&-10&1001&1&10&1\\
7&22&273&19&2&1\\
8&-39&55&4&1&1\\
9&-35&429&1&7&2\\
10&-143&42&5&1&2\\
11&-3&13&1&2&1\\
12&-26&105&1&2&1\\
13&-6&385&1&8&1\\
14&-14&65&3&2&1\\
15&-11&5&3&1&2\\
16&-1&13&-2&3&1
\end{array}
\end{equation}
All these numbers are integral. Their base discriminants and norms are
supported on $\Sigma=\{2,3,5,7,11,13\}$, so their Galois closures are
unramified outside $\Sigma$. An individual catalog field need not survive
all the cuts in $G$. We will check retention for the products used in
the construction.

For a nonnegative integer $a<2^{17}$ define
\[
 \beta(a)=\prod_{i=0}^{16}\alpha_i^{a_i},\qquad
 q_a(v)=\sum_{i=0}^{16}a_i\chi_{A_i}(v)\chi_{B_i}(v),
 \quad a=\sum_i2^ia_i.
\]
Products and signs in the second formula are taken in $\mathbb F_2$.
The norm identities determine the central extension associated to each
product.
\begin{lemma}\label{ar:kummer}
For any catalog word $a$, the field $E(\sqrt{\beta(a)})$ is Galois over
$\mathbb Q$, and its invariant Kummer class has transgression $q_a$. When this form
is nonzero, the field is a central extension of $V$ by $C_2$ with
that quadratic form. Products of the catalog radicands add their extension classes.
On the space of invariant Kummer classes in the maximal pro-$2$
$\Sigma$-unramified extension above $E$, this class map is injective.
\end{lemma}
\begin{proof}
Choose the square roots of the positive rational numbers $B_i$ compatibly
inside $E$. If $v\in V$ changes the sign of $\sqrt{A_i}$, put
\[
 u_{i,v}=\frac{z_i\sqrt{B_i}}{\alpha_i}.
\]
Otherwise put $u_{i,v}=1$. The norm identity gives
\[
 \alpha_i u_{i,v}^2=v(\alpha_i),\qquad
 u_{i,v}v(u_{i,v})=(-1)^{\chi_{A_i}(v)\chi_{B_i}(v)}.
\]
Thus $v$ lifts to an automorphism by sending
$\sqrt{\alpha_i}$ to $u_{i,v}\sqrt{\alpha_i}$. For a product radicand
use $u_v=\prod_i u_{i,v}^{a_i}$. It satisfies
\begin{equation}\label{ar:kummer-lifts}
 \beta(a)u_v^2=v(\beta(a)),\qquad
 u_vv(u_v)=(-1)^{q_a(v)}.
\end{equation}
The kernel over $E$ is central of exponent two. The squares of lifted
elements are given by $q_a$, and their commutators by its
polar form $B_{q_a}(v,w)=q_a(v+w)+q_a(v)+q_a(w)$. These squares and
commutators determine the central extension, since every word can be
collected into an ordered product of the seven lifted basis elements and
a central sign. Applying the same calculation to a product of radicands
gives the addition rule.

For injectivity let $\mathcal H$ be the subgroup of $\mathcal G$ fixing
$E$. The five-term sequence contains
\[
 H^1(\mathcal G,\mathbb F_2)\longrightarrow
 H^1(\mathcal H,\mathbb F_2)^V
 \xrightarrow{\mathrm{trg}} H^2(V,\mathbb F_2).
\]
Every quadratic character of $\mathcal G$ factors through $E$, by the
Kummer computation in Lemma~\ref{tw:complete-global-presentation}.
The first arrow is zero, so transgression is injective.
The calculation of squares and commutators above identifies this
transgression with the indicated quadratic form. In particular a
nonzero form cannot come from a square radicand in $E$, and independent
forms give independent Kummer classes.
\end{proof}

\subsection{The Completed Field and the Full Quadratic Quotient}
Define two lists of raw catalog words
\begin{align}
 \mathcal A&=(17,27184,6220,2,32768,39508,132),
       \label{ar:completed-words}\\
 \mathcal B&=(2,256,512,4096,9,65,132,2052,24576,37,32768,65536).
       \label{ar:full-words}
\end{align}
Set
\[
 N=E\bigl(\sqrt{\beta(a)}:a\in\mathcal A\bigr),\qquad
 M=E\bigl(\sqrt{\beta(b)}:b\in\mathcal B\bigr).
\]
For example, writing $\eta=\alpha_{13}\alpha_{14}$, the seven defining
radicands of $N$ are
\begin{align*}
 \gamma_1&=\alpha_0\alpha_4,&
 \gamma_2&=\alpha_4\alpha_5\alpha_9\alpha_{11}\eta,\\
 \gamma_3&=\alpha_2\alpha_3\alpha_6\alpha_{11}\alpha_{12},&
 \gamma_4&=\alpha_1,\\
 \gamma_5&=3+\sqrt{-11},&
 \gamma_6&=\alpha_2\alpha_4\alpha_6\alpha_9\alpha_{11}\alpha_{12}\gamma_5,\\
 \gamma_7&=(1+4\sqrt{-2})(19+2\sqrt{22}).
\end{align*}

These products give two subfields of the tower.
\begin{proposition}\label{ar:central-fields}
The fields $N$ and $M$ are contained in every field $K$ of
Theorem~\ref{tw:field-family}. Their degrees are
\[
 [N:\mathbb Q]=2^{14}=16384,\qquad [M:\mathbb Q]=2^{19},\qquad N\subset M.
\]
The field $M$ is exactly the field corresponding to $G/D_3G$.
\end{proposition}
\begin{proof}
Use the 28-coordinate convention dual to the restricted quadratic basis
in Section~\ref{tw:tower}: the first seven coordinates encode $v_i^2$
and the remaining 21 encode $v_iv_j$ for lexicographically ordered
$i<j$. In this convention the forms for $\mathcal A$ and $\mathcal B$
have hexadecimal masks
\begin{equation}\label{ar:central-masks}
\begin{array}{c|l}
\mathcal A&
\texttt{9014280},\ \texttt{af24500},\ \texttt{6830680},\ \texttt{b401400},\\
&\texttt{800200},\ \texttt{acf0980},\ \texttt{a138900}\\[2pt]
\mathcal B&
\texttt{b401400},\ \texttt{9140a00},\ \texttt{78c1900},\ \texttt{520e700},\\
&\texttt{2415680},\ \texttt{1c38e00},\ \texttt{a138900},\ \texttt{213900},\\
&\texttt{45f9c00},\ \texttt{1bb80},\ \texttt{800200},\ \texttt{1000}.
\end{array}
\end{equation}
To obtain these masks, expand each product of linear characters from
\eqref{ar:norm-table}, writing a diagonal term as $v_i^2$.
Binary elimination gives ranks seven and twelve, respectively, and
shows that the first span is contained in the second. Taking the binary
scalar product with each of the sixteen displayed relation rows in
Section~\ref{tw:tower} gives zero for every form in the second list.
The rows have rank 16 in dimension 28, so these twelve independent forms
span their entire annihilator.

Lemma~\ref{ar:kummer} now proves the asserted Kummer degrees and the
containment. Each extension has central kernel of exponent two over the
exponent-two group $V$, and hence class at most two and exponent at most
four. Its third Zassenhaus subgroup is thus trivial. Evaluating a defining
relator of $G$ in this group depends only on its quadratic initial, and
the higher-degree relators vanish. Since the forms annihilate every
quadratic relation, these number fields satisfy all the defining
relations of $G$, including the local cuts. Finally
$G/D_3G$ has the same order $2^{19}$ as $M$, by
Lemma~\ref{tw:new-retained-quotient}, which proves equality. Since the
fields $K$ retain $G/D_4G$, they retain both $M$ and $N$.
\end{proof}

\subsection{The Complete Representation Inventory}
For a quadratic form $q$ on $V$, write $R_q=\operatorname{rad}(B_q)$
and $g(q)=\sum_{v\in V}(-1)^{q(v)}$. The space of forms defining $N$
will be denoted $\mathcal Q=\langle q_a:a\in\mathcal A\rangle$.
To count the irreducible representations and their gamma factors, we
examine the polar forms and the action of complex conjugation.
\begin{lemma}\label{ar:representations}
The irreducible representations of $\operatorname{Gal}(N/\mathbb Q)$
comprise 128 of degree one, 256 of degree two, 456 of degree four, and
124 of degree eight. Exactly eight of the degree-four representations
are pure at infinity: four have gamma vector $(0,0,0,0)$ and four have
$(1,1,1,1)$. Every other nonlinear representation is mixed, with equal
numbers of zeros and ones in its gamma vector.
\end{lemma}
\begin{proof}
Fix a nontrivial character of $\operatorname{Gal}(N/E)$ and let
$q\in\mathcal Q\setminus\{0\}$ be its central form. Suppose
$\operatorname{rank}B_q=2m$. The corresponding central pushout has
commutator pairing $B_q$, and its center above $V$ is the inverse image
of $R_q$. There are $|R_q|=128/2^{2m}$ central characters extending the
chosen nontrivial kernel character. Each gives one irreducible
representation of dimension $2^m$.
To construct these representations, choose a maximal isotropic subspace
of $V$ for $B_q$, of dimension $7-m$ and containing $R_q$. Its inverse image is abelian.
Extend the chosen central character to it and induce to the full group.
Nondegeneracy of the induced pairing on $V/R_q$ shows, by the usual
intertwiner criterion for induction, that no element outside that
subgroup intertwines its character. The induced representation is
irreducible and has degree $2^m$. Changing the extension inside this
subgroup only conjugates the character, while changing its restriction
to the center gives distinct representations. Their squared dimensions
sum to 128, so this accounts for the entire central sector.
The trivial central character gives the 128 linear characters of $E$.

Enumeration of the 127 nonzero sums of the seven masks in
\eqref{ar:central-masks} gives
\[
\begin{array}{c|rrr}
 \operatorname{rank}B_q&2&4&6\\\hline
 \#q&8&57&62\\
 \#\{q:g(q)>0\}&5&34&48\\
 \#\{q:g(q)=0\}&3&23&14.
\end{array}
\]
The positive sums are respectively 64, 32, and 16. No negative Gauss
sum occurs. For this enumeration, evaluate each quadratic polynomial on the 128
binary vectors and compute the rank of its $7$ by $7$ polar matrix.
The resulting dimension count is
\[
 128+256\cdot2^2+456\cdot4^2+124\cdot8^2=16384.
\]

The linear map $q\mapsto B_q(c,\cdot)$ on $\mathcal Q$, where
$c=e_0$, has rank six. Its only nonzero kernel form has rank four
and is represented by $\alpha_4\alpha_6\alpha_9\alpha_{10}\alpha_{11}$.
Every form satisfies $q(c)=0$, as it annihilates the infinity square.
If $B_q(c,\cdot)\ne0$, an element conjugating $c$ to its product with
the nontrivial central sign shows that the trace of $c$ in the
representation is zero. Because $c^2=1$, its $+1$ and $-1$ eigenspaces
then have equal dimensions. For the unique remaining sector $c$ is
central in the pushout, hence acts scalarly. Its eight genus twists
split equally between the two signs. The explicit description below
identifies the even and odd twists. The dimensions of these eigenspaces
give the asserted real gamma factors.
\end{proof}

The regular representation gives the exact Artin factorization
\begin{equation}\label{ar:artin-factorization}
 \zeta_N(s)=\prod_{\rho\in\operatorname{Irr}(\operatorname{Gal}(N/\mathbb Q))}
                     L(s,\rho)^{\dim\rho}.
\end{equation}
This follows prime by prime by decomposing the regular representation
and taking the determinant on inertia invariants. We next identify every
nonlinear factor with an entire Hecke function,
which allows us to use this identity without an Artin holomorphy
assumption.

\subsection{Descent to Quadratic Hecke Characters}\label{ar:descent}
We construct the inducing characters by descending the catalog radicands
to suitable subfields of $E$. The descent gives exact identities for the
Artin functions in the analytic calculation.

If $g(q)>0$ and $\operatorname{rank}B_q=2m$, then $q$ vanishes on its
radical and the induced nondegenerate quadratic form has Arf invariant
zero. Equivalently it has a totally singular subspace $U$ of dimension
$7-m$ containing $R_q$. We may find such a subspace by enumerating the subspaces of $V$ and
testing $q(u)=0$ on each vector. Put $B=E^U$, so
$[B:\mathbb Q]=2^m$, and choose a catalog product $\beta$ representing $q$.
Let $u_v$ be the exact lifts from \eqref{ar:kummer-lifts}. On $E$ define
$T_vP=u_vv(P)$. For $v\in U$ these are commuting involutions: their
squares and commutators are respectively $(-1)^{q(v)}$ and
$(-1)^{B_q(v,w)}$.

Starting from $P=1$, take a basis $v_1,\ldots,v_{7-m}$ of $U$ and
successively replace
\begin{equation}\label{ar:projectors}
 P\longmapsto P+T_{v_j}P
 \quad\hbox{or, if this is zero,}\quad P-T_{v_j}P.
\end{equation}
At least one choice is nonzero in characteristic zero. Commutativity
preserves the eigenspaces already selected, so the final nonzero $P$
satisfies $T_vP=\pm P$ for every $v\in U$. Consequently
\[
 \eta=\beta P^2\in B,
 \qquad E(\sqrt\eta)=E(\sqrt\beta).
\]
Every operation is rational arithmetic in the basis
$\prod_i\sqrt{a_i}^{\epsilon_i}$ of $E$. At a step, if
$\eta_0=\beta P^2$, the element
$h=\beta u_vP v(P)$ satisfies
\[
 h^2=\eta_0v(\eta_0),\qquad
 \beta(P\pm T_vP)^2=\eta_0+v(\eta_0)\pm2h.
\]
At each descent step, these identities verify the new radicand, while
the nonzero test and the invariance under the processed subgroup verify
the chosen projector.

Let $L=B(\sqrt\eta)$ and let $\lambda$ be its nontrivial quadratic
Hecke character over $B$. The inverse image of $U$ in the central
pushout is the elementary abelian subgroup used in the representation
proof. Its character selecting $\sqrt\eta$ induces the required
irreducible representation. Thus
\begin{equation}\label{ar:quadratic-induction}
 L(s,\rho)=L_B(s,\lambda)=\frac{\zeta_L(s)}{\zeta_B(s)}.
\end{equation}
This is an entire nontrivial Hecke function. Its completed functional
equation has root number $+1$, because the completed Dedekind zeta
functions in the quotient have root number $+1$.
The remaining representations in this sector are obtained by multiplying
$\eta$ by rational squareclasses $D$ representing
$V^*/B_q(V)$. There are $|R_q|$ such twists. Here $B_q(V)$ denotes the
image of $v\mapsto B_q(v,\cdot)$. Twisting by a character in this image
conjugates the inducing representation, while the quotient distinguishes
its central characters.

If $g(q)=0$, choose $p\in\{5,13\}$ so that $g(q+\chi_p)>0$.
For this seven-dimensional space, exact Gauss-sum evaluation gives
26 sectors permitting either choice, nine requiring $p=5$, and five
requiring $p=13$. Thus a suitable choice exists in every sector. Put
\[
 \beta_5=-10-2\sqrt5,\qquad \beta_{13}=-26-6\sqrt{13}.
\]
Their norms are $16p$ and their central forms are $\chi_p$.
To identify these as the cyclic quartic extensions of conductor $p$,
we use their defining polynomials
\[
 f_5(X)=X^4+20X^2+80,\qquad f_{13}(X)=X^4+52X^2+208.
\]
If $r^2=\beta_p$, the conjugation rule $r\mapsto4\sqrt p/r$ has
square $r\mapsto-r$, and gives their cyclic Galois group of order four.
The following integral bases have trace-pairing discriminants $5^3$
and $13^3$, respectively:
\[
\begin{split}
 p=5:\quad &1,\quad r^2/8-r/4+3/2,\quad-r^2/8-r/4-3/2,\quad
 r^3/16-r^2/8+3r/4-1;\\
 p=13:\quad &1,\quad r^2/24-r/4+5/6,\quad-r^2/24-r/4-5/6,\quad
 r^3/48-r^2/24+11r/12-4/3.
\end{split}
\]
Integrality and the discriminants follow by substitution in $f_p$ and
rational trace-matrix arithmetic. These orders are maximal. The field is
unramified outside $p$, and its
quadratic subfield $\mathbb Q(\sqrt p)$ is ramified at $p$. Hence inertia
is the full cyclic group of order four. Tame ramification gives
discriminant exponent three, which agrees with the computed order
discriminant.
The conductor--discriminant formula gives conductor $p$ for both faithful
quartic characters, identifying them with $\xi_p$ and
$\overline{\xi_p}$ under global reciprocity.
Descend $\beta\beta_p$ by the positive construction. If $\rho'$ is
its quadratic induction and $\xi_p$ is the primitive quartic Dirichlet
character normalized by $\xi_p(2)=i$, use $\rho'\otimes\xi_p$.
The two quadratic forms add to
$(q+\chi_p)+\chi_p=q$, so this is a representation of the actual
original central field. This can also be seen directly on the two
Kummer square roots: the extra central sign from $\beta_p$ occurs
in both factors and cancels in their tensor product. Its genus twists
again give every representation in the sector. Since
$\overline{\xi_p}=\xi_p\chi_p$ and $\rho'$ is real, the genus twists
pair the complex representation with its conjugate. The resulting
functions are entire Hecke functions with root numbers of modulus one.
The bounds below do not require their phases to be known.

\subsection{The Pure Quartic Field}
The unique pure form in Lemma~\ref{ar:representations} is represented
by $\beta=\alpha_4\alpha_6\alpha_9\alpha_{10}\alpha_{11}$.
Five projector steps give
\[
 B_0=\mathbb Q(\sqrt{30},\sqrt{65}),\qquad
 \eta_0=393911-10360\sqrt{30}+23013\sqrt{65}-2400\sqrt{1950}.
\]
To verify this squareclass identity, start with $P_0=1$ and the
multipliers $u_v$ for this $\beta$. Apply
$P_j=P_{j-1}+\epsilon_jT_{v_j}P_{j-1}$ with
\[
 (v_1,\ldots,v_5)=(1,16,32,6,74),\qquad
 (\epsilon_1,\ldots,\epsilon_5)=(1,1,1,-1,-1).
\]
Exact multiplication in $E$ gives $\eta_0=\beta P_5^2/256$.
The five vectors span the simultaneous kernels of $\chi_{30}$ and
$\chi_{65}$, so their fixed field is $B_0$. All four real values are
positive: using $\sqrt{30}<6$, $\sqrt{65}<9$, and $\sqrt{1950}<45$,
each is larger than
$393911-10360\cdot6-23013\cdot9-2400\cdot45=16634$. The eight factors are
\begin{equation}\label{ar:pure-factors}
 L_D(s)=\frac{\zeta_{B_0(\sqrt{D\eta_0})}(s)}{\zeta_{B_0}(s)},
 \qquad D\in\{\pm1,\pm2,\pm3,\pm5\}.
\end{equation}
The base discriminant is
$120\cdot65\cdot312=2433600$. For every listed $D$, the maximal-order
discriminant of $B_0(\sqrt{D\eta_0})$ is
\[
 2^{14}3^4 5^4 7^2 11^2 13^4=140455850895360000.
\]
The ratio is $240240^2$, giving the conductor of each factor by the
conductor--discriminant formula. Positive $D$ have four even gamma
factors, negative $D$ four odd ones, and every root number is $+1$.
The discriminants here are those of maximal orders, as required by the
formula below.

\subsection{Exact Local Data}
The inducing characters determine the completed functions used in the
analytic estimate.
\begin{proposition}\label{ar:hecke-data}
All irreducible factors of $\zeta_N$ are the actual Dirichlet factors
and Hecke inductions just constructed. The gamma inventory is the one
in Lemma~\ref{ar:representations}. Their conductors and exceptional
Euler denominators are given by the following finite local algorithms,
with exact inducing-field, radicand, twist, and projector data in the
supplement. In particular no conjectural Artin holomorphy, conductor
identification from a truncated Euler product, or guessed root number
is used.
\end{proposition}
\begin{proof}
The preceding constructions give every irreducible representation in
the regular representation. To compute its Hecke data in a positive
sector with twist $D$, compute the maximal orders of
$B$ and $L_D=B(\sqrt{D\eta})$. The conductor-discriminant formula gives
\begin{equation}\label{ar:conductor-ratio}
 Q_\rho=\frac{|\Delta_{L_D}|}{|\Delta_B|}
       =|\Delta_B|\,N_{B/\mathbb Q}(\mathfrak d_{L_D/B}).
\end{equation}
For a prime $p$, factor $p$ in both maximal orders. With
$f(\mathfrak p/p)$ denoting residue degrees, put
\[
 D_{B,p}(T)=\prod_{\mathfrak p\mid p}(1-T^{f(\mathfrak p/p)}),
 \qquad
 D_{L_D,p}(T)=\prod_{\mathfrak P\mid p}(1-T^{f(\mathfrak P/p)}).
\]
The exact local denominator of \eqref{ar:quadratic-induction} is
\begin{equation}\label{ar:denominator-ratio}
 D_{\rho,p}(T)=D_{L_D,p}(T)/D_{B,p}(T).
\end{equation}
This quotient is a polynomial: above each prime of $B$, quadratic
splitting contributes $1-T^f$, inertia degree two contributes $1+T^f$,
and ramification contributes 1. Its degree is at most $\dim\rho$.
At a real place its character is even or odd according as $D\eta$
is positive or negative. At a complex place the induction contributes
one even and one odd real gamma factor.

We can also compute these data locally, including in the sectors with a
quartic twist. At an odd prime of $B$, first remove the
even valuation of the quadratic radicand by a square. An odd remaining
valuation is ramified with conductor exponent one. For valuation zero,
the residue quadratic character determines splitting, giving the
unramified value of the quadratic character on a uniformizer. At a
dyadic place, use the finite unit quotient and square test: a unit
$1+u$ with valuation $v(u)>2v(2)$ is a square by Hensel's lemma. Thus
squareclass and quadratic-character computations reduce to a finite
quotient of the local integer ring. Its ramification and conductor can
be determined by testing the character on the descending unit groups.
The quadratic norm criterion determines each character value. All these tests
use exact finite-ring arithmetic.

For $\rho=\operatorname{Ind}_{B}^{\mathbb Q}\lambda\otimes\xi_p$, apply
these tests to the local character
\[
 \lambda_{\mathfrak p}\,\bigl(\xi_{p,\ell}\circ
 N_{B_{\mathfrak p}/\mathbb Q_\ell}\bigr)
 \qquad(\mathfrak p\mid\ell).
\]
The local component of a primitive Dirichlet character is normalized to
have the usual Frobenius value away from its modulus. At its modulus,
its restriction to units is the inverse of the finite Dirichlet
character, and its value on the rational uniformizer is one.
For each local character, its conductor exponent is the least integer
$n$ for which it is trivial on $1+\mathfrak p^n$, with exponent zero
for an unramified character. The search is finite: its conductor is no
larger than the maximum of the two known factor conductors. Induction
then gives
\[
 a_\ell(\rho)=\sum_{\mathfrak p\mid\ell}
 f(\mathfrak p/\ell)\bigl(a_{\mathfrak p}(\lambda\xi_p\circ N)
                         +d_{B_{\mathfrak p}/\mathbb Q_\ell}\bigr),
\]
where $d_{B_{\mathfrak p}/\mathbb Q_\ell}$ is the exponent of the local
different. If a local inducing character is ramified it contributes no
Euler factor. If it is unramified, with uniformizer value $\varepsilon$,
it contributes $1-\varepsilon T^{f(\mathfrak p/\ell)}$. Multiplication over
$\mathfrak p\mid\ell$ supplies the target denominator with its actual
complex phase.

At the cancelling odd prime this calculation has a particularly short
form. Write a local uniformizer as $\pi^2=pr$, let $u$ be the normalized
quadratic radicand unit, and write $f$ for the residue degree. The
cancelled character is unramified and its value is
\begin{equation}\label{ar:cancelled-value}
 \varepsilon=(u/\mathbb F_{p^f})\,
       \xi_p((-r)^f)^{-1}.
\end{equation}
The first factor is the value of the unramified quadratic character,
and the second comes from the unit part of the norm of $\pi$. The
inverse is required by the local normalization at the modulus. This
determines each complex Euler polynomial separately.

The finite arithmetic records specify the rational defining data,
projector identities, primitive character values, and exceptional
polynomials. The supplementary checks apply the displayed identities
and local operations to these records. Independent maximal-order and
prime-factorization computations check the discriminant ratios and
local polynomials a second time. The identification of the infinite
Euler products follows from the inducing characters constructed above.
\end{proof}

\subsection{Frobenius Tests for the Prime Budget}\label{ar:frobenius}
For an ordinary prime $p\notin\Sigma$, let $v_p\in V$ be its genus
Frobenius vector. If $v_p\ne0$, a lift in either central field has square
measured by the corresponding quadratic forms. Consequently
\begin{equation}\label{ar:nonsplit-test}
 \begin{aligned}
 f_N(p)=4&\ \Longleftrightarrow\ q_a(v_p)=1\text{ for some }a\in\mathcal A,\\
 f_M(p)=4&\ \Longleftrightarrow\ q_b(v_p)=1\text{ for some }b\in\mathcal B.
 \end{aligned}
\end{equation}
If no form detects the square, the residue degree is two, since the
genus vector is nonzero. In particular a prime with $f_N(p)=2$ and
$f_M(p)=4$ has residue degree at least four in every field $K$.

If $v_p=0$, choose any square root $s_i$ of $A_i$ modulo $p$ and set
\[
 \varepsilon_i(p)=\left(\frac{x_i+y_i s_i}{p}\right)\in\{\pm1\}.
\]
The norm identity shows both that the argument is nonzero and that the
sign is independent of the choice of $s_i$: the product with its
conjugate is the nonzero square $B_i z_i^2$ modulo $p$.
For a raw word $a$, its central Frobenius sign is
$\prod_i\varepsilon_i(p)^{a_i}$. The prime splits completely in $N$
exactly when every word in \eqref{ar:completed-words} has sign $+1$.
The analogous test for $M$ uses \eqref{ar:full-words}. Any detected
central sign gives residue degree two in that field. Since $N\subset M$,
the completed-field detections are a subset of the full-field detections.
Order the seven completed predicates before the twelve full predicates,
and assign each prime to its first nonzero predicate. This gives
disjoint sets of primes detected in the two fields, which we will use
for the finite prime corrections in the analytic estimate.
